\section[title={2024-09-30 - Der Umgang mit dem
Vater},reference={2024-09-30-Der Umgang mit dem Vater}]

\startframedtask
{\bf Erläutern} Sie, wie die Figuren der Mutter (S. 100ff.), der Tochter
(S. 45f. + S. 99) und des Sohns (S. 45f. + S. 48 + S. 115) mit den
bisherigen Konflikten umgehen.

\stopframedtask

\startsectionlevel[title={Mutter},reference={mutter}]

\startitemize[packed]
\item
  Trinkt Alkohol (\quotation{eine Falsche Spätlese} S. 100 Unten)
\item
  nimmt die gesammte \quotation{Schuld} auf sich (\quotation{alles ist
  meine Schuld} S. 100)
\item
  Singt und spiel Klavier
\item
  Mutter argumentiert systematisch / logisch
  \startitemize[packed]
  \item
    Rechnet dem Vater die Rechnungen vor, das sie es sich nicht leisten
    können (S. 101)
  \item
    Vater ignoriert sie, gibt großzügig Geld auf seinen Reisen aus.
  \stopitemize
\stopitemize

\stopsectionlevel

\startsectionlevel[title={Tochter},reference={tochter}]

\startitemize[packed]
\item
  Raucht Zigaretten
\item
  Verdient sich ihr eigenes Taschengeld
\item
  Lügt um eigenen Vorteiel zu haben
\stopitemize

\startframedtask
Erläutern Sie, was müsste die Familie ggf. mit Therapie lernen, um
Konflikte künftig konstruktiv bewältigen zu können.

\stopframedtask

\stopsectionlevel

\startsectionlevel[title={DDR},reference={ddr}]

\startitemize[packed]
\item
  Mauerbau 13. August 1961
\stopitemize

\stopsectionlevel

\startsectionlevel[title={Sil:},reference={sil}]

\startframedtask
\ldots{}

\stopframedtask

\stopsectionlevel
